\documentclass[11pt,a4paper]{article}
\usepackage[T1]{fontenc}\usepackage[utf8]{inputenc}
\usepackage{lmodern,microtype,amsmath,amssymb,amsthm,mathtools}
\usepackage[margin=24mm]{geometry}
\usepackage{booktabs,array}\usepackage[hyphens]{url}\usepackage[hidelinks]{hyperref}
\allowdisplaybreaks\setlength{\emergencystretch}{3em}
\newtheorem{theorem}{Theorem}[section]\newtheorem{lemma}[theorem]{Lemma}
\newtheorem{proposition}[theorem]{Proposition}
\newcommand{\Li}{L_{-8}(2)}\newcommand{\hh}{H^{(3)}}
\hypersetup{pdftitle={An existing-WZ reduction of Sun Conjecture 5.39},pdfauthor={Reiwa}}
\title{An existing-WZ reduction of Sun Conjecture 5.39\\[5pt]\large A global proof candidate, not a new-method claim}
\author{Reiwa --- Tournament round 8}\date{16 November 2026}
\begin{document}\maketitle
\begin{abstract}
A one-dimensional specialization of Hou and Sun's existing two-parameter
Wilf--Zeilberger pair reduces Sun's Conjecture 5.39 to a previously proved
cubic harmonic Ramanujan-type identity and an elementary logarithmic
integral. We give the rational certificate, regularized parameter family,
uniform boundary estimates, and coefficient extraction. This is a proposed
complete derivation of the original assertion, submitted for independent
verification. It is not a claim of an entirely new method: the main WZ seed
and the necessary modular evaluation are prior results, so the tournament's
substantive-new-method gate G5 is explicitly not passed. Priority and the
absence of an earlier identical specialization are not certified.
\end{abstract}

\section{Original assertion and necessary prior results}
Write $H_n^{(r)}=\sum_{j=1}^n j^{-r}$, $H_0^{(r)}=0$,
$G=\sum_{j\ge0}(-1)^j/(2j+1)^2$, and $P(n)=28n^2+10n+1$.
Sun's Conjecture 5.39, dated 2025-02-04, is equation (5.38) on page 32 of
\cite[v10]{Sun}. It states exactly
\begin{equation}\label{target}
 \mathcal S:=\sum_{n=0}^{\infty}w_n
 \left\{P(n)(28H_{2n}^{(3)}-3H_n^{(3)})+\frac8{2n+1}\right\}
 =\frac{75\zeta(3)}\pi-24G,
\end{equation}
where
\begin{equation}\label{weight}
 w_n=\frac{\binom{2n}{n}^{5}}
 {(6n+1)(-64)^n\binom{3n}{n}\binom{6n}{3n}}
 =\left(-\frac1{27}\right)^n
 \frac{(1/2)_n^5}{(5/6)_n(7/6)_n(n!)^3}.
\end{equation}
In particular the $n=0$ summand is $8$. The final equality in
\eqref{weight} follows from the gamma multiplication formulas, or from
initial values and consecutive-term ratios. There are no restrictions on a
prime or any finite truncation in this assertion.

We use the following established evaluations, not new conjectural inputs:
\begin{align}
 v_n&=\left(-\frac18\right)^n\frac{(1/2)_n^3}{(n!)^3}(6n+1),
 &\sum_{n\ge0}v_n&=\frac{2\sqrt2}{\pi},\label{known1}\\
 \sum_{n\ge0}v_n\left(H_{2n}^{(3)}-\frac7{64}H_n^{(3)}\right)
 &=\frac{57\zeta(3)}{16\sqrt2\pi}-\Li.\label{known2}
\end{align}
Here $\Li=\sum_{j\ge1}\chi_{-8}(j)/j^2$, with $\chi_{-8}$ zero on even
integers and taking values $1,1,-1,-1$ on residues $1,3,5,7$ modulo $8$.
Equation \eqref{known1} is Ramanujan's standard identity, also used in
\cite[p.~17]{HS}. Equation \eqref{known2} is exactly
\cite[p.~5, (1.30)]{SZ}; that paper attributes an earlier proof to Li and Chu
and supplies a modular proof through Theorem 1.2, Table 2 and (3.39).
The term $-\Li$ in \eqref{known2} is essential; it must not be dropped.

\section{Explicit parameter family and rational certificate}
For $|x|<1/16$ and integers $n,k\ge0$, define
\begin{equation}\label{F}
 F_x(n,k)=\frac{6(-1)^{n+1}4^k}{\pi^2}
 \frac{\Gamma(n-2x+\tfrac12)^2\Gamma(n+k+x+\tfrac12)^2
 \Gamma(n+4x+\tfrac12)\Gamma(n+k-2x+\tfrac12)}
 {\Gamma(k+\tfrac32)\Gamma(n-2x)\Gamma(n+x)^2
 \Gamma(3n+2k+\tfrac32)}.
\end{equation}
Reciprocals of gamma functions are entire; in particular $F_0(0,k)=0$.
This family is obtained from the prior WZ seed
\cite[p.~17, (2.1)]{HS} by replacing its $n$ by $n+x$, setting $a=3x,c=0$,
and removing the scalar phase $e^{-6\pi i x}$ after reflection of
$\Gamma(-1/2-k)$. The explicitly written real-sign family \eqref{F}, rather
than an ambiguous complex power $(-1)^x$, is used throughout.

Put $D=(n-2x)(n+x)^2$, and
\begin{align}
 Q_x(n,k)=-\frac18\bigl[&(24n+4)k^2+(72n^2+36n+4)k\notag\\
 &+56n^3+48n^2+12n+1-(24n+12)x^2+32x^3\bigr].\label{Q}
\end{align}
Initially away from $D=0$, set
\begin{equation}\label{R}
 R_x(n,k)=\frac{(2k+1)Q_x(n,k)}{D(4k+6n+3)},\qquad
 G_x(n,k)=R_x(n,k)F_x(n,k).
\end{equation}
At a zero of $D$, define $G_x$ by cancellation: replace
$\Gamma(n-2x)D\Gamma(n+x)^2$ in the product by
$\Gamma(n-2x+1)\Gamma(n+x+1)^2$. This gives a holomorphic function on
the entire disk, including $(n,x)=(0,0)$.

\begin{lemma}\label{cert}
These functions satisfy
\begin{equation}\label{WZ}
 F_x(n+1,k)-F_x(n,k)=G_x(n,k+1)-G_x(n,k).
\end{equation}
\end{lemma}
\begin{proof}
For generic $x$ let $r_n=F_x(n+1,k)/F_x(n,k)$ and
$r_k=F_x(n,k+1)/F_x(n,k)$. The gamma recurrence gives
\begin{align*}
 r_n&=-\frac{(n-2x+\tfrac12)^2(n+k+x+\tfrac12)^2
 (n+4x+\tfrac12)(n+k-2x+\tfrac12)}
 {D(3n+2k+\tfrac32)(3n+2k+\tfrac52)(3n+2k+\tfrac72)},\\
 r_k&=\frac{4(n+k+x+\tfrac12)^2(n+k-2x+\tfrac12)}
 {(k+\tfrac32)(3n+2k+\tfrac32)(3n+2k+\tfrac52)}.
\end{align*}
The finite polynomial check $r_kR_x(n,k+1)-R_x(n,k)=r_n-1$
proves the claim. For a completely denominator-free certificate, put
$A=2k+2n+2x+1$, $B=2k+2n-4x+1$, $C=4k+6n+3$, and
$N=(2n-4x+1)^2A^2(2n+8x+1)B$. The check is precisely
\[
 32A^2B Q_x(n,k+1)-8(2k+1)Q_x(n,k)(C+2)(C+4)
 +N+8DC(C+2)(C+4)=0.
\]
Substitution of \eqref{Q} proves this by expansion. Holomorphic
continuation gives \eqref{WZ} also at the removable exceptional points.
\end{proof}

Here are the two relevant boundary expressions, with all constants fixed:
\begin{align}
 A_x(n)&=\frac{(-1)^n}{2\pi^{5/2}}
 \frac{\Gamma(n-2x+\tfrac12)^3\Gamma(n+x+\tfrac12)^2
 \Gamma(n+4x+\tfrac12)}
 {\Gamma(n-2x+1)\Gamma(n+x+1)^2\Gamma(3n+\tfrac32)},\label{Ax}\\
 G_x(n,0)&=A_x(n)\left(P(n)-12x^2+\frac{32x^3}{2n+1}\right),\label{G0}\\
 g_x(n):=\lim_{k\to\infty}G_x(n,k)
 &=\frac{3(-1)^n(6n+1)}{2^{3n+3/2}\pi^{3/2}}
 \frac{\Gamma(n-2x+\tfrac12)^2\Gamma(n+4x+\tfrac12)}
 {\Gamma(n-2x+1)\Gamma(n+x+1)^2}.\label{g}
\end{align}
Indeed $A_0(n)=w_n$, $g_0(n)=3\sqrt2\,v_n/4$, and
\[
 -8Q_x(n,0)=(2n+1)(P(n)-12x^2)+32x^3,\qquad
 \lim_{k\to\infty}\frac{R_x(n,k)}{k^2}=-\frac{6n+1}{4D}.
\]
The duplication formula and uniform gamma-ratio asymptotics yield
\eqref{g} directly from \eqref{F}; no phase convention is needed.

\section{All boundary limits and coefficient interchanges}
\begin{lemma}\label{normal}
On every closed disk $|x|\le\rho<1/16$, all three series in
\begin{equation}\label{master}
 \sum_{n\ge0}G_x(n,0)=\sum_{n\ge0}g_x(n)+\sum_{k\ge0}F_x(0,k)
\end{equation}
converge normally. Moreover \eqref{master} holds throughout the disk.
Consequently its Taylor coefficients may be extracted term by term.
\end{lemma}
\begin{proof}
All numerator gamma arguments have positive real part on the disk.
The fixed-denominator cancellation defining $G_x$ removes its apparent
pole at $x=0,n=0$. The standard estimate
$\Gamma(y+u)/\Gamma(y)=y^u(1+O(1/y))$, uniform for $u$ in a compact set,
shows that products of such ratios whose shifts sum to zero are bounded
uniformly for real $y\ge1/2$.

For $n\ge1$, divide \eqref{F} by $F_0(n,k)$ and group the shifts at
arguments $n+1/2$, $n+k+1/2$, and $n$. The respective shifts have sums
$-2-2+4=0$, $1+1-2=0$, and $-2+1+1=0$.
Thus $|F_x(n,k)|\le C_\rho|F_0(n,k)|$, uniformly in $n,k$.
Set
\[
 C_n=\frac6{2^{3n+1/2}\pi^{3/2}}
       \frac{\Gamma(n+1/2)^3}{\Gamma(n)^3}=O(8^{-n}n^{3/2}).
\]
Duplication gives
\[
 |F_0(n,k)|=C_n
 \frac{\Gamma(k+n+1/2)^3}
 {\Gamma(k+3/2)\Gamma(k+3n/2+3/4)\Gamma(k+3n/2+5/4)}.
\]
Use the gamma recurrence once in the first and third denominator factors.
The three remaining arguments
$k+1/2$, $k+3n/2+3/4$, $k+3n/2+1/4$ have mean $k+n+1/2$.
Log-convexity therefore implies
\[
 |F_0(n,k)|\le\frac{C_n}{(k+1/2)(k+3n/2+1/4)}
 \le\frac{C_n}{(k+1/2)^2}\quad(n\ge1).
\]
In particular $\sum_k|F_x(n,k)|\le C_\rho C_n\pi^2/2\to0$
uniformly as $n\to\infty$. This verifies the usually delicate moving
boundary, not merely the limit for each fixed $k$.

For the row $n=0$, duplication and the same balanced gamma-ratio estimate
give $F_x(0,k)=O_\rho((k+1)^{-2})$, including the entire factor
$1/(\Gamma(-2x)\Gamma(x)^2)$. The two remaining series have, uniformly,
$G_x(n,0)=O_\rho((1+n^2)27^{-n})$ and
$g_x(n)=O_\rho((1+n)8^{-n})$. These bounds follow from \eqref{Ax}--\eqref{g},
balanced shifts, and \eqref{weight}; looser polynomial bounds suffice.
Thus all required series are normally convergent.

Sum \eqref{WZ} first over $0\le n<N$, $0\le k<K$. Let $K\to\infty$
with $N$ fixed, using the row convergence and \eqref{g}. Then let
$N\to\infty$, using the uniform moving-boundary estimate. The result is
\eqref{master}, with its displayed signs. Normal convergence of
holomorphic functions, or Cauchy's coefficient formula on a smaller disk,
justifies all subsequent coefficient extractions.
\end{proof}

\section{The cubic coefficient and its elementary boundary}
Use $[x^3]$ for an ordinary Taylor coefficient, not the third derivative.
The first logarithmic derivatives of both $A_x(n)$ and $g_x(n)$ vanish
at zero because the corresponding sums of gamma shifts are zero.
Recall
\[
 \psi''(n+1/2)=-14\zeta(3)+16H_{2n}^{(3)}-2H_n^{(3)},\qquad
 \psi''(n+1)=-2\zeta(3)+2H_n^{(3)}.
\]
The sums of cubes of numerator shifts are $42$ for $A_x$ and $48$ for
$g_x$; the denominator sum is $-6$ in both cases. Consequently
\begin{align}
 [x^3]A_x(n)&=w_n\{4(28H_{2n}^{(3)}-3H_n^{(3)})-100\zeta(3)\},\label{Acub}\\
 [x^3]g_x(n)&=g_0(n)\{128H_{2n}^{(3)}-14H_n^{(3)}-114\zeta(3)\}.\label{gcub}
\end{align}
For example, the logarithmic cubic coefficient in \eqref{Acub} is
$7\psi''(n+1/2)+\psi''(n+1)$. There are no products of first and second
logarithmic derivatives, since the first one is zero. The term $-12x^2$
in \eqref{G0} similarly contributes nothing to the cubic coefficient.
Equations \eqref{G0}, \eqref{Acub} give
\begin{equation}\label{leftcub}
 [x^3]\sum_{n\ge0}G_x(n,0)=4\mathcal S-100\zeta(3)\sum_{n\ge0}w_nP(n).
\end{equation}
At $x=0$, \eqref{master} and \eqref{known1} give
$\sum_nw_nP(n)=\sum_ng_0(n)=3/\pi$.
From \eqref{known1}, \eqref{known2} and \eqref{gcub}, we also obtain
\begin{equation}\label{rightcub}
 [x^3]\sum_{n\ge0}g_x(n)=-96\sqrt2\,\Li.
\end{equation}

It remains to account for the $n=0$ row; it cannot be discarded merely
because its constant coefficient vanishes. Since
$1/(\Gamma(-2x)\Gamma(x)^2)=-2x^3+O(x^4)$, direct substitution into
\eqref{F} yields
\begin{equation}\label{fcub}
 f_k:=[x^3]F_x(0,k)=\frac{12\,4^k}{\sqrt\pi}
 \frac{\Gamma(k+1/2)^3}{\Gamma(k+3/2)\Gamma(2k+3/2)}.
\end{equation}
\begin{lemma}\label{elementary}
$\displaystyle\sum_{k\ge0}f_k=96(\sqrt2\,\Li-G)$.
\end{lemma}
\begin{proof}
The beta integral rewrites \eqref{fcub} as
\[
 f_k=\frac{24\binom{2k}{k}}{2k+1}
       \int_0^1t^{k-1/2}(1-t)^k\,dt.
\]
Tonelli's theorem applies since all terms are nonnegative. The generating
function $\sum_{k\ge0}\binom{2k}{k}u^k/(2k+1)
=\arcsin(2\sqrt u)/(2\sqrt u)$ holds on $0\le u\le1/4$, including the
endpoint by monotone convergence. Therefore
\begin{align*}
 \sum_kf_k
 &=12\int_0^1\frac{\arcsin(2\sqrt{t(1-t)})}{t\sqrt{1-t}}\,dt\\
 &=48\left\{\int_0^{\pi/4}\frac{\theta}{\sin\theta}\,d\theta
 +\int_{\pi/4}^{\pi/2}\frac{\pi/2-\theta}{\sin\theta}\,d\theta\right\}\\
 &=96\left\{\int_{\pi/8}^{\pi/4}\log\tan u\,du
                  -\int_0^{\pi/8}\log\tan u\,du\right\}.
\end{align*}
Here $t=\sin^2\theta$ and
$\arcsin(\sin2\theta)=2\min(\theta,\pi/2-\theta)$.
Integration by parts uses $(\log\tan(\theta/2))'=\csc\theta$;
the endpoint terms cancel and $\theta\log\theta\to0$ at zero.
The integrals are absolutely convergent. Finally the integrated Fourier
series of $\log\tan u$, obtainable by Abel limiting, gives
\[
 \int_0^a\log\tan u\,du=-\sum_{\substack{j\ge1\\j\text{ odd}}}
                   \frac{\sin(2ja)}{j^2},\quad 0<a\le\pi/4.
\]
It yields $-G$ at $a=\pi/4$ and $-\Li/\sqrt2$ at $a=\pi/8$.
The resulting absolutely convergent sine series justify both evaluations,
and the claimed value follows.
\end{proof}

\begin{theorem}[Proposed full derivation of the original assertion]
The established identities \eqref{known1}--\eqref{known2}, together with
the verified certificate \eqref{WZ}, give \eqref{target}.
\end{theorem}
\begin{proof}
Take $[x^3]$ of \eqref{master}, justified by Lemma~\ref{normal}.
Equations \eqref{leftcub}, \eqref{rightcub} and Lemma~\ref{elementary} give
\[
 4\mathcal S-\frac{300\zeta(3)}\pi
 =-96\sqrt2\,\Li+96(\sqrt2\,\Li-G)=-96G.
\]
Division by four proves exactly \eqref{target}.
\end{proof}

\section{Independent checks, literature audit, and honest gate verdict}
\paragraph{Finite exact checks.}
The accompanying \texttt{verify.py} checks the denominator-free WZ
identity symbolically, the $k=0$ polynomial, and the large-$k$ normalization.
It checks the literal binomial-to-Pochhammer conversion and the harmonic
coefficient identity for $0\le n\le104$ using exact rational arithmetic.
It also encloses both sides of \eqref{target} independently. Indeed
$|w_n|\le27^{-n}$ and, for $n\ge1$, the bracket in \eqref{target} is at most
$1646n^2$, since $H_n^{(3)}<3/2$ and $P(n)\le39n^2$.
Thus the tail from $n=N$ is bounded by
\[
 1646\,27^{-N}\left\{\frac{N^2}{1-q}
 +\frac{2Nq}{(1-q)^2}+\frac{q(1+q)}{(1-q)^3}\right\},\qquad q=1/27.
\]
For $N=105$ this is less than $9.602\cdot10^{-144}$.
Exact rational enclosures of $\pi,\zeta(3),G$ use Machin's arctangent
formula and positive Euler transforms of their alternating series.
The recorded interval for left side minus right side lies inside
$[-1.178\cdot10^{-134},7.428\cdot10^{-135}]$.
The common value begins
\[
 6.71382044208451930745100115968774161836919465852487180883129035354349711713.
\]
These tests are reproducible supporting evidence, not a replacement for
Sections 2--4. The auxiliary numerical script checks the two boundary
coefficients, including the nonzero $\Li$ term, separately.

\paragraph{Prior-art scope and version discipline.}
The official arXiv records were checked on 2026-11-16. The original source
is still v10 (6 February 2026), and still labels the precise statement a
conjecture. Hou--Sun is v1 (16 April 2026); its Theorem 1.2 states the
harmonic-free and first-harmonic identities on this same weight, not
\eqref{target}. Its seed (2.1) is the direct antecedent of \eqref{F}.
Sun--Zhou v2 (2 March 2026) establishes the separate cubic harmonic input
\eqref{known2}; it does not state \eqref{target}. Au v4
(9 November 2026) explicitly develops coefficient extraction from WZ
seeds and discusses the interaction with modular proofs. This is serious
method-level prior art, not evidence of a new general technique here.

Targeted searches for the conjecture number, the polynomial
$28n^2+10n+1$, the cubic harmonic combination, and the constants
$75\zeta(3)/\pi-24G$ located no explicit earlier proof of the exact target
in the inspected results. This is only a bounded search. No assertion
that all published, unpublished, equivalent, or translated proofs were
excluded is made. In particular an earlier implicit consequence of the
same seed and modular identity is entirely possible.

The supplied Verifiable Mathematical Reasoning Corpus was audited at
commit \texttt{5887cdea4898f0a3aa5bf7a8e0829697bd2b0f4b}, covering all
190 available TEX files. The saved audit records the literal source,
formula, label, and WZ-method matches: zero direct hits for the eight target/source
patterns, generic WZ hits in five files, and coefficient-extraction/differentiation
hits in 26 files. A lexical scan does not establish
mathematical inequivalence, and the direct prior WZ seed is acknowledged
regardless of whether it occurs in that corpus.

\begin{center}\small
\begin{tabular}{p{.08\linewidth}p{.83\linewidth}}
\toprule
Gate&Verdict for this submission\\\midrule
G1&Pass: a named, dated existing conjecture, not a newly invented target.\\
G2&Pass: the complete original single identity is \eqref{target}.\\
G3&Provisional only: latest original labels it open; no exact subsequent proof
located in the bounded audit. Global unresolved status is not certified.\\
G4&Complete proof candidate given, with unconditional established inputs and
all parameter/boundary steps specified; independent review still pending.\\
G5&\textbf{Not passed.} This adapts the existing Hou--Sun WZ seed and uses an
existing modular evaluation. No substantive-new-method or tournament-win claim.\\
G6&Bounded audit completed; method overlap is explicit. Historical priority,
implicit-equivalence exclusion, and exhaustive nonduplication are not certified.\\
\bottomrule
\end{tabular}
\end{center}
\noindent\textbf{Submission status.} Preserve this as an independently
checkable proof candidate and an explicit failure of the new-method gate,
not as an announced new solution satisfying all six gates.

\begin{thebibliography}{9}
\bibitem{Sun} Z.-W. Sun, \emph{New series involving binomial coefficients (II)},
arXiv:2307.03086v10, 6 February 2026, p.~32, Conjecture 5.39 and (5.38).
\url{https://arxiv.org/abs/2307.03086v10}.
\bibitem{HS} Q.-H. Hou and Z.-W. Sun, \emph{Evaluations of some series via the
WZ method}, arXiv:2604.15172v1, 16 April 2026, Theorem 1.2 and p.~17, (2.1).
\url{https://arxiv.org/abs/2604.15172v1}.
\bibitem{SZ} Z.-W. Sun and Y. Zhou, \emph{Series involving central binomial
coefficients and higher-order harmonic numbers}, arXiv:2602.12091v2,
2 March 2026, (1.30), Theorem 1.2, Table 2, and (3.39).
\url{https://arxiv.org/abs/2602.12091v2}.
\bibitem{Au} K. C. Au, \emph{Discovering hypergeometric series with harmonic
numbers via Wilf--Zeilberger seeds}, arXiv:2602.08721v4, 9 November 2026.
\url{https://arxiv.org/abs/2602.08721v4}.
\bibitem{Corpus} QEEDDD, \emph{Verifiable Mathematical Reasoning Corpus},
inspected 190-file snapshot, commit recorded above.
\url{https://github.com/QEEDDD/Verifiable-Mathematical-Reasoning-Corpus}.
\end{thebibliography}
\end{document}
